% themestats.tex: measure a theme for the website's live slide.
%
% % themestats.tex: measure a theme for the website's live slide.
%
% % themestats.tex: measure a theme for the website's live slide.
%
% % themestats.tex: measure a theme for the website's live slide.
%
% \input{themestats} just before \end{document} in demo.tex. It adds one
% slide at the very end (so no other page number moves) and writes the same
% numbers to <jobname>-stats.txt. Leave it out of the decks you publish by
% building them without \def\ShowStats.
%
% Everything is in TeX points (pt), rounded to 2 decimals.
% On the website, 1pt = 100cqw / paper width.

% \ifdefined\ShowStats

% ---- Number formatting (expl3, part of every LaTeX since 2020) -----------
\ExplSyntaxOn
% A length as a number with exactly two decimals: 455.24408pt -> 455.24
\cs_new:Npn \tfs_num:n #1
  { \exp_args:Ne \__tfs_fixed:n { \fp_eval:n { round ( \dim_to_fp:n {#1} * 100 ) } } }
\cs_new:Npn \__tfs_fixed:n #1
  {
    \int_div_truncate:nn {#1} {100} .
    \int_compare:nNnT { \int_mod:nn {#1} {100} } < {10} {0}
    \int_mod:nn {#1} {100}
  }
% Pad text with spaces to a fixed width: on the left (numbers line up on the
% right) or on the right (labels, so the next column starts in one place).
\cs_new:Npn \__tfs_lpad:nn #1#2
  { \prg_replicate:nn { \int_max:nn {0} { #1 - \str_count:n {#2} } } {~} #2 }
\cs_new:Npn \__tfs_rpad:nn #1#2
  { #2 \prg_replicate:nn { \int_max:nn {0} { #1 - \str_count:n {#2} } } {~} }
\cs_generate_variant:Nn \__tfs_lpad:nn { ne }
\cs_generate_variant:Nn \__tfs_rpad:nn { ne }
% Names usable outside expl3 syntax
\cs_new:Npn \TFnum #1 { \tfs_num:n {#1} }
\cs_new:Npn \TFlpad #1#2 { \__tfs_lpad:ne {#1} {#2} }
\cs_new:Npn \TFrpad #1#2 { \__tfs_rpad:ne {#1} {#2} }
\ExplSyntaxOff

\makeatletter

% ---- Output: one line in the text file + one table row on the slide ------
% Column widths in the text file, in characters.
\def\tf@wl{16}   % label
\def\tf@wn{10}   % each number column
\def\tf@wb{14}   % the baselineskip column (its heading is longer)
\newwrite\tf@out
\immediate\openout\tf@out=\jobname-stats.txt

% \tf@line{label}{value}: a text line, label padded to one width.
\newcommand\tf@line[2]{\immediate\write\tf@out{\TFrpad{\tf@wl}{#1}#2}}

% A length: label, then the number right-aligned in its column.
\newcommand\tf@len[2]{% #1 = label, #2 = length
  \tf@line{#1}{\TFlpad{\tf@wn}{\TFnum{#2}}}%
  \texttt{#1} & \TFnum{#2} &&\\
}

% Height of a beamer template, as drawn on this slide.
\newcommand\tf@tpl[2]{% #1 = label, #2 = template name
  \setbox0=\vbox{\usebeamertemplate*{#2}}%
  \xdef\tf@h{\the\dimexpr\ht0+\dp0\relax}% keep it: each table cell is its own group
  \tf@len{#1}{\tf@h}%
}

% A font: size, 1ex (x-height) and baselineskip, in three columns.
\newcommand\tf@font[2]{% #1 = label, #2 = font switch
  \begingroup #2% switch font inside a group, keep the numbers with \xdef
    \xdef\tf@size{\f@size pt}%
    \xdef\tf@ex{\the\fontdimen5\font}%
    \xdef\tf@bl{\the\baselineskip}%
  \endgroup
  \tf@line{#1}{\TFlpad{\tf@wn}{\TFnum{\tf@size}}%
               \TFlpad{\tf@wn}{\TFnum{\tf@ex}}%
               \TFlpad{\tf@wb}{\TFnum{\tf@bl}}}%
  \texttt{#1} & \TFnum{\tf@size} & \TFnum{\tf@ex} & \TFnum{\tf@bl}\\
}

% A heading row for the font columns.
\newcommand\tf@fonthead{%
  \noalign{% \noalign must come first in a table row, so the writes go inside it
    \immediate\write\tf@out{}%
    \tf@line{font}{\TFlpad{\tf@wn}{size}\TFlpad{\tf@wn}{1ex}\TFlpad{\tf@wb}{baselineskip}}%
    \smallskip}%
  \textit{font} & \textit{size} & \textit{1ex} & \textit{baselineskip}\\
}

\begin{frame}[noframenumbering]{Theme stats: \ThemeName}{Measured with a one-line subtitle}
  \tiny
  \tf@line{theme}{\ThemeName\space[\ThemeOptions]}%
  \immediate\write\tf@out{}%
  \begin{tabular}{@{}lrrr@{}}
    \multicolumn{4}{@{}l}{\texttt{theme}\quad \ThemeName\space[\ThemeOptions]}\\
    \tf@len{paper width}{\paperwidth}
    \tf@len{paper height}{\paperheight}
    \tf@len{text width}{\textwidth}
    \tf@len{left margin}{\beamer@leftmargin}
    \tf@len{right margin}{\beamer@rightmargin}
    \tf@tpl{headline}{headline}
    \tf@tpl{frametitle}{frametitle}
    \tf@tpl{footline}{footline}
    \tf@fonthead
    \tf@font{normalsize}{\normalsize}
    \tf@font{scriptsize}{\scriptsize}
    \tf@font{footnotesize}{\footnotesize}
    \tf@font{small}{\small}
    \tf@font{large}{\large}
    \tf@font{Large}{\Large}
    \tf@font{frametitle}{\usebeamerfont{frametitle}}
    \tf@font{framesubtitle}{\usebeamerfont{framesubtitle}}
    \tf@font{title}{\usebeamerfont{title}}
    \tf@font{block title}{\usebeamerfont{block title}}
    \tf@font{footline}{\usebeamerfont{footline}}
  \end{tabular}
\end{frame}

\immediate\closeout\tf@out
\makeatother
% \fi
 just before \end{document} in demo.tex. It adds one
% slide at the very end (so no other page number moves) and writes the same
% numbers to <jobname>-stats.txt. Leave it out of the decks you publish by
% building them without \def\ShowStats.
%
% Everything is in TeX points (pt), rounded to 2 decimals.
% On the website, 1pt = 100cqw / paper width.

% \ifdefined\ShowStats

% ---- Number formatting (expl3, part of every LaTeX since 2020) -----------
\ExplSyntaxOn
% A length as a number with exactly two decimals: 455.24408pt -> 455.24
\cs_new:Npn \tfs_num:n #1
  { \exp_args:Ne \__tfs_fixed:n { \fp_eval:n { round ( \dim_to_fp:n {#1} * 100 ) } } }
\cs_new:Npn \__tfs_fixed:n #1
  {
    \int_div_truncate:nn {#1} {100} .
    \int_compare:nNnT { \int_mod:nn {#1} {100} } < {10} {0}
    \int_mod:nn {#1} {100}
  }
% Pad text with spaces to a fixed width: on the left (numbers line up on the
% right) or on the right (labels, so the next column starts in one place).
\cs_new:Npn \__tfs_lpad:nn #1#2
  { \prg_replicate:nn { \int_max:nn {0} { #1 - \str_count:n {#2} } } {~} #2 }
\cs_new:Npn \__tfs_rpad:nn #1#2
  { #2 \prg_replicate:nn { \int_max:nn {0} { #1 - \str_count:n {#2} } } {~} }
\cs_generate_variant:Nn \__tfs_lpad:nn { ne }
\cs_generate_variant:Nn \__tfs_rpad:nn { ne }
% Names usable outside expl3 syntax
\cs_new:Npn \TFnum #1 { \tfs_num:n {#1} }
\cs_new:Npn \TFlpad #1#2 { \__tfs_lpad:ne {#1} {#2} }
\cs_new:Npn \TFrpad #1#2 { \__tfs_rpad:ne {#1} {#2} }
\ExplSyntaxOff

\makeatletter

% ---- Output: one line in the text file + one table row on the slide ------
% Column widths in the text file, in characters.
\def\tf@wl{16}   % label
\def\tf@wn{10}   % each number column
\def\tf@wb{14}   % the baselineskip column (its heading is longer)
\newwrite\tf@out
\immediate\openout\tf@out=\jobname-stats.txt

% \tf@line{label}{value}: a text line, label padded to one width.
\newcommand\tf@line[2]{\immediate\write\tf@out{\TFrpad{\tf@wl}{#1}#2}}

% A length: label, then the number right-aligned in its column.
\newcommand\tf@len[2]{% #1 = label, #2 = length
  \tf@line{#1}{\TFlpad{\tf@wn}{\TFnum{#2}}}%
  \texttt{#1} & \TFnum{#2} &&\\
}

% Height of a beamer template, as drawn on this slide.
\newcommand\tf@tpl[2]{% #1 = label, #2 = template name
  \setbox0=\vbox{\usebeamertemplate*{#2}}%
  \xdef\tf@h{\the\dimexpr\ht0+\dp0\relax}% keep it: each table cell is its own group
  \tf@len{#1}{\tf@h}%
}

% A font: size, 1ex (x-height) and baselineskip, in three columns.
\newcommand\tf@font[2]{% #1 = label, #2 = font switch
  \begingroup #2% switch font inside a group, keep the numbers with \xdef
    \xdef\tf@size{\f@size pt}%
    \xdef\tf@ex{\the\fontdimen5\font}%
    \xdef\tf@bl{\the\baselineskip}%
  \endgroup
  \tf@line{#1}{\TFlpad{\tf@wn}{\TFnum{\tf@size}}%
               \TFlpad{\tf@wn}{\TFnum{\tf@ex}}%
               \TFlpad{\tf@wb}{\TFnum{\tf@bl}}}%
  \texttt{#1} & \TFnum{\tf@size} & \TFnum{\tf@ex} & \TFnum{\tf@bl}\\
}

% A heading row for the font columns.
\newcommand\tf@fonthead{%
  \noalign{% \noalign must come first in a table row, so the writes go inside it
    \immediate\write\tf@out{}%
    \tf@line{font}{\TFlpad{\tf@wn}{size}\TFlpad{\tf@wn}{1ex}\TFlpad{\tf@wb}{baselineskip}}%
    \smallskip}%
  \textit{font} & \textit{size} & \textit{1ex} & \textit{baselineskip}\\
}

\begin{frame}[noframenumbering]{Theme stats: \ThemeName}{Measured with a one-line subtitle}
  \tiny
  \tf@line{theme}{\ThemeName\space[\ThemeOptions]}%
  \immediate\write\tf@out{}%
  \begin{tabular}{@{}lrrr@{}}
    \multicolumn{4}{@{}l}{\texttt{theme}\quad \ThemeName\space[\ThemeOptions]}\\
    \tf@len{paper width}{\paperwidth}
    \tf@len{paper height}{\paperheight}
    \tf@len{text width}{\textwidth}
    \tf@len{left margin}{\beamer@leftmargin}
    \tf@len{right margin}{\beamer@rightmargin}
    \tf@tpl{headline}{headline}
    \tf@tpl{frametitle}{frametitle}
    \tf@tpl{footline}{footline}
    \tf@fonthead
    \tf@font{normalsize}{\normalsize}
    \tf@font{scriptsize}{\scriptsize}
    \tf@font{footnotesize}{\footnotesize}
    \tf@font{small}{\small}
    \tf@font{large}{\large}
    \tf@font{Large}{\Large}
    \tf@font{frametitle}{\usebeamerfont{frametitle}}
    \tf@font{framesubtitle}{\usebeamerfont{framesubtitle}}
    \tf@font{title}{\usebeamerfont{title}}
    \tf@font{block title}{\usebeamerfont{block title}}
    \tf@font{footline}{\usebeamerfont{footline}}
  \end{tabular}
\end{frame}

\immediate\closeout\tf@out
\makeatother
% \fi
 just before \end{document} in demo.tex. It adds one
% slide at the very end (so no other page number moves) and writes the same
% numbers to <jobname>-stats.txt. Leave it out of the decks you publish by
% building them without \def\ShowStats.
%
% Everything is in TeX points (pt), rounded to 2 decimals.
% On the website, 1pt = 100cqw / paper width.

% \ifdefined\ShowStats

% ---- Number formatting (expl3, part of every LaTeX since 2020) -----------
\ExplSyntaxOn
% A length as a number with exactly two decimals: 455.24408pt -> 455.24
\cs_new:Npn \tfs_num:n #1
  { \exp_args:Ne \__tfs_fixed:n { \fp_eval:n { round ( \dim_to_fp:n {#1} * 100 ) } } }
\cs_new:Npn \__tfs_fixed:n #1
  {
    \int_div_truncate:nn {#1} {100} .
    \int_compare:nNnT { \int_mod:nn {#1} {100} } < {10} {0}
    \int_mod:nn {#1} {100}
  }
% Pad text with spaces to a fixed width: on the left (numbers line up on the
% right) or on the right (labels, so the next column starts in one place).
\cs_new:Npn \__tfs_lpad:nn #1#2
  { \prg_replicate:nn { \int_max:nn {0} { #1 - \str_count:n {#2} } } {~} #2 }
\cs_new:Npn \__tfs_rpad:nn #1#2
  { #2 \prg_replicate:nn { \int_max:nn {0} { #1 - \str_count:n {#2} } } {~} }
\cs_generate_variant:Nn \__tfs_lpad:nn { ne }
\cs_generate_variant:Nn \__tfs_rpad:nn { ne }
% Names usable outside expl3 syntax
\cs_new:Npn \TFnum #1 { \tfs_num:n {#1} }
\cs_new:Npn \TFlpad #1#2 { \__tfs_lpad:ne {#1} {#2} }
\cs_new:Npn \TFrpad #1#2 { \__tfs_rpad:ne {#1} {#2} }
\ExplSyntaxOff

\makeatletter

% ---- Output: one line in the text file + one table row on the slide ------
% Column widths in the text file, in characters.
\def\tf@wl{16}   % label
\def\tf@wn{10}   % each number column
\def\tf@wb{14}   % the baselineskip column (its heading is longer)
\newwrite\tf@out
\immediate\openout\tf@out=\jobname-stats.txt

% \tf@line{label}{value}: a text line, label padded to one width.
\newcommand\tf@line[2]{\immediate\write\tf@out{\TFrpad{\tf@wl}{#1}#2}}

% A length: label, then the number right-aligned in its column.
\newcommand\tf@len[2]{% #1 = label, #2 = length
  \tf@line{#1}{\TFlpad{\tf@wn}{\TFnum{#2}}}%
  \texttt{#1} & \TFnum{#2} &&\\
}

% Height of a beamer template, as drawn on this slide.
\newcommand\tf@tpl[2]{% #1 = label, #2 = template name
  \setbox0=\vbox{\usebeamertemplate*{#2}}%
  \xdef\tf@h{\the\dimexpr\ht0+\dp0\relax}% keep it: each table cell is its own group
  \tf@len{#1}{\tf@h}%
}

% A font: size, 1ex (x-height) and baselineskip, in three columns.
\newcommand\tf@font[2]{% #1 = label, #2 = font switch
  \begingroup #2% switch font inside a group, keep the numbers with \xdef
    \xdef\tf@size{\f@size pt}%
    \xdef\tf@ex{\the\fontdimen5\font}%
    \xdef\tf@bl{\the\baselineskip}%
  \endgroup
  \tf@line{#1}{\TFlpad{\tf@wn}{\TFnum{\tf@size}}%
               \TFlpad{\tf@wn}{\TFnum{\tf@ex}}%
               \TFlpad{\tf@wb}{\TFnum{\tf@bl}}}%
  \texttt{#1} & \TFnum{\tf@size} & \TFnum{\tf@ex} & \TFnum{\tf@bl}\\
}

% A heading row for the font columns.
\newcommand\tf@fonthead{%
  \noalign{% \noalign must come first in a table row, so the writes go inside it
    \immediate\write\tf@out{}%
    \tf@line{font}{\TFlpad{\tf@wn}{size}\TFlpad{\tf@wn}{1ex}\TFlpad{\tf@wb}{baselineskip}}%
    \smallskip}%
  \textit{font} & \textit{size} & \textit{1ex} & \textit{baselineskip}\\
}

\begin{frame}[noframenumbering]{Theme stats: \ThemeName}{Measured with a one-line subtitle}
  \tiny
  \tf@line{theme}{\ThemeName\space[\ThemeOptions]}%
  \immediate\write\tf@out{}%
  \begin{tabular}{@{}lrrr@{}}
    \multicolumn{4}{@{}l}{\texttt{theme}\quad \ThemeName\space[\ThemeOptions]}\\
    \tf@len{paper width}{\paperwidth}
    \tf@len{paper height}{\paperheight}
    \tf@len{text width}{\textwidth}
    \tf@len{left margin}{\beamer@leftmargin}
    \tf@len{right margin}{\beamer@rightmargin}
    \tf@tpl{headline}{headline}
    \tf@tpl{frametitle}{frametitle}
    \tf@tpl{footline}{footline}
    \tf@fonthead
    \tf@font{normalsize}{\normalsize}
    \tf@font{scriptsize}{\scriptsize}
    \tf@font{footnotesize}{\footnotesize}
    \tf@font{small}{\small}
    \tf@font{large}{\large}
    \tf@font{Large}{\Large}
    \tf@font{frametitle}{\usebeamerfont{frametitle}}
    \tf@font{framesubtitle}{\usebeamerfont{framesubtitle}}
    \tf@font{title}{\usebeamerfont{title}}
    \tf@font{block title}{\usebeamerfont{block title}}
    \tf@font{footline}{\usebeamerfont{footline}}
  \end{tabular}
\end{frame}

\immediate\closeout\tf@out
\makeatother
% \fi
 just before \end{document} in demo.tex. It adds one
% slide at the very end (so no other page number moves) and writes the same
% numbers to <jobname>-stats.txt. Leave it out of the decks you publish by
% building them without \def\ShowStats.
%
% Everything is in TeX points (pt), rounded to 2 decimals.
% On the website, 1pt = 100cqw / paper width.

% \ifdefined\ShowStats

% ---- Number formatting (expl3, part of every LaTeX since 2020) -----------
\ExplSyntaxOn
% A length as a number with exactly two decimals: 455.24408pt -> 455.24
\cs_new:Npn \tfs_num:n #1
  { \exp_args:Ne \__tfs_fixed:n { \fp_eval:n { round ( \dim_to_fp:n {#1} * 100 ) } } }
\cs_new:Npn \__tfs_fixed:n #1
  {
    \int_div_truncate:nn {#1} {100} .
    \int_compare:nNnT { \int_mod:nn {#1} {100} } < {10} {0}
    \int_mod:nn {#1} {100}
  }
% Pad text with spaces to a fixed width: on the left (numbers line up on the
% right) or on the right (labels, so the next column starts in one place).
\cs_new:Npn \__tfs_lpad:nn #1#2
  { \prg_replicate:nn { \int_max:nn {0} { #1 - \str_count:n {#2} } } {~} #2 }
\cs_new:Npn \__tfs_rpad:nn #1#2
  { #2 \prg_replicate:nn { \int_max:nn {0} { #1 - \str_count:n {#2} } } {~} }
\cs_generate_variant:Nn \__tfs_lpad:nn { ne }
\cs_generate_variant:Nn \__tfs_rpad:nn { ne }
% Names usable outside expl3 syntax
\cs_new:Npn \TFnum #1 { \tfs_num:n {#1} }
\cs_new:Npn \TFlpad #1#2 { \__tfs_lpad:ne {#1} {#2} }
\cs_new:Npn \TFrpad #1#2 { \__tfs_rpad:ne {#1} {#2} }
\ExplSyntaxOff

\makeatletter

% ---- Output: one line in the text file + one table row on the slide ------
% Column widths in the text file, in characters.
\def\tf@wl{16}   % label
\def\tf@wn{10}   % each number column
\def\tf@wb{14}   % the baselineskip column (its heading is longer)
\newwrite\tf@out
\immediate\openout\tf@out=\jobname-stats.txt

% \tf@line{label}{value}: a text line, label padded to one width.
\newcommand\tf@line[2]{\immediate\write\tf@out{\TFrpad{\tf@wl}{#1}#2}}

% A length: label, then the number right-aligned in its column.
\newcommand\tf@len[2]{% #1 = label, #2 = length
  \tf@line{#1}{\TFlpad{\tf@wn}{\TFnum{#2}}}%
  \texttt{#1} & \TFnum{#2} &&\\
}

% Height of a beamer template, as drawn on this slide.
\newcommand\tf@tpl[2]{% #1 = label, #2 = template name
  \setbox0=\vbox{\usebeamertemplate*{#2}}%
  \xdef\tf@h{\the\dimexpr\ht0+\dp0\relax}% keep it: each table cell is its own group
  \tf@len{#1}{\tf@h}%
}

% A font: size, 1ex (x-height) and baselineskip, in three columns.
\newcommand\tf@font[2]{% #1 = label, #2 = font switch
  \begingroup #2% switch font inside a group, keep the numbers with \xdef
    \xdef\tf@size{\f@size pt}%
    \xdef\tf@ex{\the\fontdimen5\font}%
    \xdef\tf@bl{\the\baselineskip}%
  \endgroup
  \tf@line{#1}{\TFlpad{\tf@wn}{\TFnum{\tf@size}}%
               \TFlpad{\tf@wn}{\TFnum{\tf@ex}}%
               \TFlpad{\tf@wb}{\TFnum{\tf@bl}}}%
  \texttt{#1} & \TFnum{\tf@size} & \TFnum{\tf@ex} & \TFnum{\tf@bl}\\
}

% A heading row for the font columns.
\newcommand\tf@fonthead{%
  \noalign{% \noalign must come first in a table row, so the writes go inside it
    \immediate\write\tf@out{}%
    \tf@line{font}{\TFlpad{\tf@wn}{size}\TFlpad{\tf@wn}{1ex}\TFlpad{\tf@wb}{baselineskip}}%
    \smallskip}%
  \textit{font} & \textit{size} & \textit{1ex} & \textit{baselineskip}\\
}

\begin{frame}[noframenumbering]{Theme stats: \ThemeName}{Measured with a one-line subtitle}
  \tiny
  \tf@line{theme}{\ThemeName\space[\ThemeOptions]}%
  \immediate\write\tf@out{}%
  \begin{tabular}{@{}lrrr@{}}
    \multicolumn{4}{@{}l}{\texttt{theme}\quad \ThemeName\space[\ThemeOptions]}\\
    \tf@len{paper width}{\paperwidth}
    \tf@len{paper height}{\paperheight}
    \tf@len{text width}{\textwidth}
    \tf@len{left margin}{\beamer@leftmargin}
    \tf@len{right margin}{\beamer@rightmargin}
    \tf@tpl{headline}{headline}
    \tf@tpl{frametitle}{frametitle}
    \tf@tpl{footline}{footline}
    \tf@fonthead
    \tf@font{normalsize}{\normalsize}
    \tf@font{scriptsize}{\scriptsize}
    \tf@font{footnotesize}{\footnotesize}
    \tf@font{small}{\small}
    \tf@font{large}{\large}
    \tf@font{Large}{\Large}
    \tf@font{frametitle}{\usebeamerfont{frametitle}}
    \tf@font{framesubtitle}{\usebeamerfont{framesubtitle}}
    \tf@font{title}{\usebeamerfont{title}}
    \tf@font{block title}{\usebeamerfont{block title}}
    \tf@font{footline}{\usebeamerfont{footline}}
  \end{tabular}
\end{frame}

\immediate\closeout\tf@out
\makeatother
% \fi
